%%%%%%%%%%%%%%%%%%%%%%%%%%%%%%%%%%%%%%%%%%%%%%%%%%%%%%%%%%%%%%%%%%%%%%
% Lớp 10 | Véctơ (chưa xét tọa độ) | Tổng và hiệu của hai véctơ
% Dạng 0H5?2-4: Tính độ dài của véctơ tổng, hiệu
% Số câu: 17  (sinh tự động bởi tach_cau.py; từ nay thầy sửa trực tiếp file này)
%%%%%%%%%%%%%%%%%%%%%%%%%%%%%%%%%%%%%%%%%%%%%%%%%%%%%%%%%%%%%%%%%%%%%%

%%%%%%%%%%%%%%%% TRẮC NGHIỆM 4 PHƯƠNG ÁN %%%%%%%%%%%%%%%%

% nguồn: MaHoaNH_TracNghiem_Lop10
\begin{ex}%[0H5N2-4]%[ma:26991fc5]
	Cho tam giác đều $ABC$ có cạnh bằng $a$. Khi đó $\left|\overrightarrow{CB}-\overrightarrow{CA}\right|$  bằng
	 \choice
		{$a\sqrt{3}$} 
		{$a$} 
		 {$2a$}
		{\True  $\dfrac{a\sqrt{3}}{2}$}
	 \loigiai{
	Ta có $\left|\overrightarrow{CB}-\overrightarrow{CA}\right|=\left|\overrightarrow{AB}\right|=a$  do tam giác đều $ABC$ có cạnh bằng $a$.}
\end{ex}

% nguồn: MaHoaNH_TracNghiem_Lop10
\begin{ex}%[0H5N2-4]%[ma:5250c152]
	Cho tam giác $ABC$ đều cạnh $a$. Độ dài $\overrightarrow{AB}+\overrightarrow{BC}$  bằng
	 \choice
		{\True  $a$}
		 {$\dfrac{a\sqrt{3}}{2}$}
		 {$2a$}
		{$a\sqrt{3}$} 
	\loigiai{ 
		Ta có 
	$\left|\overrightarrow{AB}+\overrightarrow{BC}\right|=\left|\overrightarrow{AC}\right|=AC=a$.} 
\end{ex}

% nguồn: MaHoaNH_TracNghiem_Lop10
\begin{ex}%[0H5H2-4]%[ma:5ad07d82]
	Cho hình vuông  $ABCD$ cạnh $a$. Tính $\left|\overrightarrow{AB}-\overrightarrow{DA}\right|$ là
	 \choice
		{$\left|\overrightarrow{AB}-\overrightarrow{DA}\right|=2a$} 
		{$\left|\overrightarrow{AB}-\overrightarrow{DA}\right|=0$} 
		 {\True $\left|\overrightarrow{AB}-\overrightarrow{DA}\right|=a\sqrt{2}$}
		 {$\left|\overrightarrow{AB}-\overrightarrow{DA}\right|=a$}
	 \loigiai{
	 \immini{
		Ta  có $\overrightarrow{AB}=\overrightarrow{DC}$ nên $$\left|\overrightarrow{AB}-\overrightarrow{DA}\right|=\left|\overrightarrow{DC}-\overrightarrow{DA}\right|=\left|\overrightarrow{AC}\right|=AC=a\sqrt{2}.$$
		}{\begin{tikzpicture}[line join=round,line cap=round,>=stealth,scale=1,font=\footnotesize]
		\foreach \x/\y/\n in {0/0/A,2/0/D,0/2/B} \coordinate (\n) at  (\x,\y);
		\coordinate  (C) at ($(B)+(D)-(A)$);
		\coordinate  (O) at ($(B)!0.5!(D)$);
		\coordinate (M) at  ($(A)!1/2!(B)$);
		 \draw (A)--(B)--(C)--(D)--(A)--(C)
			 (D)--(B);
		\foreach \t/\g  in {A/-135,D/-45,C/45,B/135,O/-90}{
			 \draw[fill=black] (\t) circle (1pt) node[shift={(\g:8pt)}]{$ \t $};
		 }
	\end{tikzpicture}
		} 
	 }
\end{ex}

% nguồn: MaHoaNH_TracNghiem_Lop10
\begin{ex}%[0H5H2-4]%[ma:54f9a74a]
	Cho hình chữ nhật $ABCD$, $AB=3$,  $AD=4$. Tính $\left|\overrightarrow{AB}+\overrightarrow{AD}\right|$.
	 \choice
		{$\left|\overrightarrow{AB}+\overrightarrow{AD}\right|=8$} 
		 {$\left|\overrightarrow{AB}+\overrightarrow{AD}\right|=7$}
		 {$\left|\overrightarrow{AB}+\overrightarrow{AD}\right|=6$}
		{\True $\left|\overrightarrow{AB}+\overrightarrow{AD}\right|=5$} 
	\loigiai{ 
	 \immini{
		Ta có $\overrightarrow{AB}+\overrightarrow{AD}=\overrightarrow{AC}$.\\ 
		Suy ra  $\left|\overrightarrow{AB}+\overrightarrow{AD}\right|=\left|\overrightarrow{AC}\right|=AC=\sqrt{A{B^2}+B{C^2}}=\sqrt{3^2+4^2}=5$.
		}{\begin{tikzpicture}[line join=round,line cap=round,>=stealth,scale=1,font=\footnotesize]
		\foreach  \x/\y/\n in {0/0/A,2/0/D,0/1.5/B} \coordinate (\n) at (\x,\y);
		\coordinate (C)  at ($(B)+(D)-(A)$);
		\coordinate (O)  at ($(B)!0.5!(D)$);
		\coordinate (M) at  ($(A)!1/2!(B)$);
		 \draw (A)--(B)--(C)--(D)--(A)--(C)
			(D)--(B); 
		\foreach \t/\g in  {A/-135,D/-45,C/45,B/135,O/-90}{
			\draw[fill=black] (\t) circle (1pt) node[shift={(\g:8pt)}]{$ \t  $};
		} 
	\end{tikzpicture}
		 }
	} 
\end{ex}

% nguồn: MaHoaNH_TracNghiem_Lop10
\begin{ex}%[0H5H2-4]%[ma:f6618000]
	Cho hình vuông $ABCD$ cạnh $a$, tâm $O$. Tính độ  dài của véc-tơ $\overrightarrow{OA}+\overrightarrow{OB}$.
	 \choice
		{$2a$} 
		{$\dfrac{a}{2}$} 
		{\True  $a$}
		 {$3a$}
	 \loigiai{
	\immini{ 
		Gọi $M$ là trung điểm của $AB$.\\ 
		 Suy ra $\left|\overrightarrow{OA}+\overrightarrow{OB}\right|=2\left|\overrightarrow{OM}\right|=2OM=AD=a$.
		}{\begin{tikzpicture}[line join=round,line cap=round,>=stealth,scale=1,font=\footnotesize]
		\foreach \x/\y/\n in {0/0/A,2/0/B,0/-2/D} \coordinate (\n) at  (\x,\y);
		\coordinate (C)  at ($(B)+(D)-(A)$);
		 \coordinate (O) at ($(B)!0.5!(D)$);
		\coordinate (M) at  ($(A)!1/2!(B)$);
		\draw (A)--(B)--(C)--(D)--(A)--(C) 
			(D)--(B) 
			(M)--(O); 
		\foreach  \t/\g in {A/135,D/-135,C/-45,B/45,O/-90,M/90}{
			\draw[fill=black] (\t) circle (1pt) node[shift={(\g:8pt)}]{$ \t $}; 
		 }
	\end{tikzpicture}
		} 
	 }
\end{ex}

% nguồn: MaHoaNH_TracNghiem_Lop10
\begin{ex}%[0H5H2-4]%[ma:2af8b5d4]
	 Cho hình vuông $ABCD$ cạnh $a$. Tính $S=\left| 2\overrightarrow{AD}+\overrightarrow{DB}\right|$.
	\choice 
		{$S=a$} 
		 {$S=a\sqrt{3}$}
		{\True $S=a\sqrt{2}$} 
		 {$S=a\sqrt{5}$}
	 \loigiai{
	\immini{ 
		Ta có $2\overrightarrow{AD}+\overrightarrow{DB}=\overrightarrow{AD}+\overrightarrow{AB}=\overrightarrow{AC}$.\\ 
		 Suy ra $\left| 2\overrightarrow{AD}+\overrightarrow{DB}\right|=\left|\overrightarrow{AC}\right|=AC=a\sqrt{2}$.\\
		Vậy  $S=a\sqrt{2}$.
		}{\begin{tikzpicture}[line join=round,line cap=round,>=stealth,scale=1,font=\footnotesize]
		\foreach \x/\y/\n in {0/0/A,2/0/D,0/2/B} \coordinate (\n) at (\x,\y); 
		\coordinate  (C) at ($(B)+(D)-(A)$);
		 \coordinate (O) at ($(B)!0.5!(D)$);
		\coordinate (M)  at ($(A)!1/2!(B)$);
		\draw  (A)--(B)--(C)--(D)--(A)--(C)
			(D)--(B); 
		\foreach  \t/\g in {A/-135,D/-45,C/45,B/135,O/-90}{
			 \draw[fill=black] (\t) circle (1pt) node[shift={(\g:8pt)}]{$ \t $};
		} 
	\end{tikzpicture}
		} 
	 }
\end{ex}

% nguồn: MaHoaNH_TracNghiem_Lop10
\begin{ex}%[0H5V2-4]%[ma:961d12fa]
	Cho hình bình hành $ABCD$ có $AB=a$, $AB\perp BD$, $\widehat{BAD}=60^\circ $. Gọi  $E$, $F$ lần lượt là trung điểm của $BD$, $AD$. Độ dài véc-tơ $\overrightarrow{BE}+\overrightarrow{AF}$ là
	 \choice
		{\True $\dfrac{a\sqrt{13}}{2}$} 
		{$\dfrac{a\sqrt{10}}{2}$} 
		{$\dfrac{a\sqrt{7}}{2}$} 
		{$ 2a$} 
	 \loigiai{
	\immini{ 
		Gọi $H$ là  trung điểm $AB$.\\
		Ta có  $\overrightarrow{BE}=\overrightarrow{ED}$, $\overrightarrow{AF}=\overrightarrow{HE}$ nên suy ra
		$$\overrightarrow{BE}+\overrightarrow{AF}=\overrightarrow{HE}+\overrightarrow{ED}=\overrightarrow{HD}$$ 
		Suy  ra $\left|\overrightarrow{BE}+\overrightarrow{AF}\right|=\left|\overrightarrow{HD}\right|=HD$.\\
		Ta lại có $BD=AB\tan 60^\circ=a\sqrt{3}$.\\ $HD=\sqrt{HB^2+BD^2}=\sqrt{\left(\dfrac{a}{2}\right)^2+\left(a\sqrt{3}\right)^2}=\dfrac{a\sqrt{13}}{2}$.\\ 
		Vậy  $\left|\overrightarrow{BE}+\overrightarrow{AF}\right|=\dfrac{a\sqrt{13}}{2}$.
		}{\begin{tikzpicture}[line join=round,line cap=round,>=stealth,scale=1,font=\footnotesize]
		\foreach \x/\y/\n in {0/0/A,4/0/D} \coordinate (\n) at  (\x,\y);
		\coordinate (B) at ($(A)+({2*cos (60)},{2*sin (60)})$); 
		 \coordinate (C) at ($(B)+(D)-(A)$);
		\coordinate  (E) at ($(D)!1/2!(B)$);
		 \coordinate (F) at ($(A)!1/2!(D)$);
		\coordinate (H) at  ($(A)!1/2!(B)$);
		 \draw (A)--(B)--(C)--(D)--(A)
			(B)--(D) 
			 (E)--(F)
			 (D)--(H)--(E);
		\foreach \t/\g in  {A/225,B/135,C/45,D/-45,E/60,F/-90,H/180}{
			\draw[fill=black] (\t) circle  (1pt) node[shift={(\g:8pt)}]{$ \t $};
		 }
		\draw pic[thin,draw,angle radius=2mm]  {right angle= A--B--D};
	\end{tikzpicture}
		} 
	 }
\end{ex}


%%%%%%%%%%%%%%%% ĐÚNG/SAI %%%%%%%%%%%%%%%%

% nguồn: MaHoaNH_DungSai_Lop10
\begin{ex}%[0H5H2-4]%[ma:f5d7fa21]
	Cho  hình thoi $ABCD$ cạnh $a$, có $\widehat{BAD}=60^{\circ}$. Gọi $O$ là giao điểm hai đường chéo.
	\choiceTF 
		 {\True $\triangle BAD$, $\triangle BCD$ đều cạnh $a$}
		{$\left|\overrightarrow{AB}+\overrightarrow{AD}\right|=a\sqrt{2}$} 
		{\True $\left|\overrightarrow{BA}-\overrightarrow{BC}\right|=a\sqrt{3}$} 
		 {\True $\left|\overrightarrow{OB}-\overrightarrow{DC}\right|=\dfrac{a\sqrt{3}}{2}$}
	\loigiai{ 
		\begin{center}
			\begin{tikzpicture}[line join=round,line cap=round,>=stealth,scale=1,font=\footnotesize]
					\foreach \x/\y/\n in {0/0/A}  \coordinate (\n) at (\x,\y);
					\coordinate (D) at ($(A)+({2*cos (30)},{2*sin  (30)})$);
					\coordinate  (B) at ($(A)+({2*cos (-30)},{2*sin (-30)})$);
					\coordinate  (C) at ($(D)+(B)-(A)$);
					\coordinate (O) at ($(A)!0.5!(C)$); 
					\draw  (C)--(A)--(B)--(C)--(D)--(A)
						 (B)--(D);
					\foreach \t/\g in  {A/180,B/-30,C/0,D/90,O/-45}{
						\draw[fill=black] (\t) circle (1pt)  node[shift={(\g:10pt)}]{$ \t $};
					 }

					\draw pic[thin,draw,angle radius=2mm] {right angle=  A--O--D};
				\end{tikzpicture}
		\end{center}
		\begin{itemchoice}
			\itemch Vì giả thiết cho ta $\triangle BAD$,  $\triangle BCD$ đều cạnh $a$.
			\itemch Vì $\left|\overrightarrow{AB}+\overrightarrow{AD}\right|=\left|\overrightarrow{AC}\right|=AC=2AO=2\sqrt{AB^2-BO^2}=a\sqrt{3}$. 
			 \itemch Vì $\left|\overrightarrow{BA}-\overrightarrow{BC}\right|=\left|\overrightarrow{CA}\right|=CA=a\sqrt{3}$.
			\itemch Vì  $\left|\overrightarrow{OB}-\overrightarrow{DC}\right|=\left|\overrightarrow{DO}-\overrightarrow{DC}\right|=\left|\overrightarrow{CO}\right|=CO=AO=\dfrac{a\sqrt{3}}{2}$.
		\end{itemchoice}
	 }
\end{ex}

% nguồn: MaHoaNH_DungSai_Lop10
\begin{ex}%[0H5V2-4]%[ma:6c9ecf66]
	Cho hình vuông  $ABCD$ có cạnh bằng $a$. Khi đó
	 \choiceTF
		 {\True $\overrightarrow{AD}+\overrightarrow{AB}=\overrightarrow{AC}$}
		 {\True Gọi $E$ là điểm đối xứng với $B$ qua $C$. Khi đó $ADEC$ là hình thang}
		 {\True $\left|\overrightarrow{AD}+\overrightarrow{AB}\right|=a\sqrt{2}$}
		 {$\left|\overrightarrow{AD}+\overrightarrow{AC}\right|=a\sqrt{3}$}
	\loigiai{ 
		\begin{itemchoice}
			\itemch Vì theo quy tắc hình bình hành ta có  $\overrightarrow{AD}+\overrightarrow{AB}=\overrightarrow{AC}$.
			\itemch Vì với $E$ là điểm đối xứng với  $B$ qua $C$, ta có $CE=AD$ và $CE\parallel AD$.
		\immini{Suy ra $ADEC$ là hình bình hành, hay $ADEC$ cũng là hình  thang.\\
			}{\begin{tikzpicture}[line join=round,line cap=round,>=stealth,scale=.7,font=\footnotesize]
		\foreach \x/\y/\n in {0/0/A,0/3/B,3/0/D}  \coordinate (\n) at (\x,\y);
		\coordinate  (C) at ($(B)+(D)-(A)$);
		\coordinate (E) at ($(C)+(C)-(B)$); 
		\draw (A)--(B)--(C)--(D)--(A)--(C) 
			 (B)--(D)--(E)--(C);
		 \foreach \t/\g in {A/-135,C/45,B/135,D/-45,E/45}{
			 \draw[fill=black] (\t) circle (1pt) node[shift={(\g:8pt)}]{$ \t $};
		 }
	\end{tikzpicture}}
			\itemch  Vì theo định lí Pytago ta có $AC=\sqrt{AB^2+BC^2}=\sqrt{a^2+a^2}=a\sqrt{2}$.\\
			Vậy $\left|\overrightarrow{AD}+\overrightarrow{AB}\right|=\left|\overrightarrow{AC}\right|=AC=a\sqrt{2}$. 
			 \itemch Vì ta có $\overrightarrow{AD}+\overrightarrow{AC}=\overrightarrow{AE}$. Áp dụng định lí Pytago cho tam giác $ABE$, ta có\\
			 $AE^2=AB^2+BE^2=a^2+\left(2a\right)^2=5a^2\Rightarrow AE=a\sqrt{5}$.\\
			Vậy $\left|\overrightarrow{AD}+\overrightarrow{AC}\right|=AE=a\sqrt{5}$. 
		\end{itemchoice}
	} 
\end{ex}


%%%%%%%%%%%%%%%% TRẢ LỜI NGẮN %%%%%%%%%%%%%%%%

% nguồn: MaHoaNH_TraLoiNgan_Lop10
\begin{bt}%[0H5H2-4]%[ma:39c96f04]
	Cho hình chữ nhật $ABCD$, $AB=3$, $AD=4$.  Tính $\left|\overrightarrow{AB}+\overrightarrow{AD}\right|$.
	 \shortans{5}
	\loigiai{ 
	Ta có  $\overrightarrow{AB}+\overrightarrow{AD}=\overrightarrow{AC}$ $\Rightarrow\left|\overrightarrow{AB}+\overrightarrow{AD}\right|=\left|\overrightarrow{AC}\right|=AC=\sqrt{AB^2+BC^2}=5$.}
\end{bt}

% nguồn: MaHoaNH_TraLoiNgan_Lop10
\begin{bt}%[0H5H2-4]%[ma:453ef310]
	Cho hình vuông $ABCD$  có độ dài cạnh bằng $2$, $M$ là trung điểm $BC$. Tính $\left|\overrightarrow{AB}+\overrightarrow{BM}\right|$ (kết quả làm tròn đến hàng phần trăm).
	\shortans{2,24} 
	\loigiai{ 
	 \immini{
		 Ta có $\left|\overrightarrow{AB}+\overrightarrow{BM}\right|=\left|\overrightarrow{AM}\right|=AM$.\\
		Theo định lí  Py-ta-go, ta có $ AM=\sqrt{AB^2+BM^2}=\sqrt{2^2+1^2}=\sqrt{5}$.\\
		Vậy $\left|\overrightarrow{AB}+\overrightarrow{BM}\right|=AM=\sqrt{5}\approx  2{,}24$.
		}{\begin{tikzpicture}[line join=round,line cap=round,>=stealth,scale=.7,font=\footnotesize]
		\foreach \x/\y/\n in {0/0/A,0/3/B,3/0/D} \coordinate (\n)  at (\x,\y);
		\coordinate (C) at  ($(B)+(D)-(A)$);
		\coordinate (M)  at ($(C)!0.5!(B)$);
		\draw (A)--(B)--(C)--(D)--(A)--(M); 
		 \foreach \t/\g in {A/-135,C/45,B/135,D/-45,M/90}{
			\draw[fill=black] (\t)  circle (1pt) node[shift={(\g:8pt)}]{$ \t $};
		 }
	\end{tikzpicture}
		 }
	} 
\end{bt}

% nguồn: MaHoaNH_TraLoiNgan_Lop10
\begin{bt}%[0H5H2-4]%[ma:f079471e]
	Cho tam giác vuông  $ABC$ có các cạnh góc vuông là $AB=1$, $AC=2$.
	Điểm $M$ thỏa mãn $\overrightarrow{AC}-\overrightarrow{AB}=\overrightarrow{AM}$. Tính độ dài véc-tơ $\overrightarrow{AM}$ (kết quả là  tròn đến hàng phần chục).
	\shortans{2,2} 
	 \loigiai{
	 \immini{
		 Ta có $\overrightarrow{AC}-\overrightarrow{AB}=\overrightarrow{AM}\Leftrightarrow \overrightarrow{BC}=\overrightarrow{AM}$.\\
		Suy  ra $ABCM$ là hình bình hành.\\
		 Ta có $\left|\overrightarrow{AM}\right|=AM=BC=\sqrt{1^2+2^2}=\sqrt{5}\approx 2{,}2$.
		}{\begin{tikzpicture}[line join=round,line cap=round,>=stealth,scale=1,font=\footnotesize]
		\foreach \x/\y/\n in {0/0/A,2/0/C,0/1/B} \coordinate (\n) at (\x,\y); 
		\coordinate  (M) at ($(A)+(C)-(B)$);
		\draw (A)--(B)--(C)--(M)--(A)--(C); 
		\foreach \t/\g in {A/-135,C/45,M/-45,B/135}{ 
			\draw[fill=black] (\t) circle (1pt) node[shift={(\g:8pt)}]{$  \t $};
		 }
		\draw pic[thin,draw,angle radius=2mm]  {right angle= B--A--C};
	\end{tikzpicture}
		} 
	} 
\end{bt}

% nguồn: MaHoaNH_TraLoiNgan_Lop10
\begin{bt}%[0H5H2-4]%[ma:86830678]
	Cho hình thang $ABCD$ có hai đáy $AB=1$, $CD=2$. Gọi $M$, $N$ lần lượt là trung điểm  $AD$ và $BC$. Tính $\left|\overrightarrow{DM}-\overrightarrow{BA}-\overrightarrow{CN}\right|$.
	 \shortans{1,5}
	 \loigiai{
	 \immini{ Ta có
		$$\begin{aligned}
		\left|\overrightarrow{DM}-\overrightarrow{BA}-\overrightarrow{CN}\right|&=\left|\overrightarrow{MA}+\overrightarrow{AB}+\overrightarrow{NC}\right|= \left|\overrightarrow{MA}+\overrightarrow{AB}+\overrightarrow{BN}\right|\\ 
		&=\left|\overrightarrow{MN}\right|=MN=\dfrac{AB+CD}{2}=\dfrac{3a}{2}. 
	\end{aligned}$$
	}{\begin{tikzpicture}[line join=round,line cap=round,>=stealth,scale=.5,font=\footnotesize]
		\foreach \x/\y/\n in {0/0/A,3/0/B,-1/-3/D,5/-3/C} \coordinate  (\n) at (\x,\y);
		\path 
			($(A)!1/2!(D)$) coordinate  (M)
			($(C)!1/2!(B)$) coordinate (N); 
		 \draw[->,thick] (D)--(M);
		 \draw[->,thick] (B)--(A);
		 \draw[->,thick] (M)--(N);
		\draw[->,thick]  (C)--(N);
		 \draw (A)--(B)--(C)--(D)--(A);
		\foreach \t/\g in {A/180,B/0,C/0,D/180,M/150,N/40}{ 
			\draw[fill=black] (\t) circle (1pt) node[shift={(\g:10pt)}]{$  \t $};
		 }
	\end{tikzpicture}
	 }
 }
\end{bt}

% nguồn: MaHoaNH_TraLoiNgan_Lop10
\begin{bt}%[0H5V2-4]%[ma:57e16d90]
	Cho hình vuông $ABCD$ có  độ dài cạnh bằng $2\sqrt{2}$. Tính $\left|\overrightarrow{AB}+\overrightarrow{AC}+\overrightarrow{AD}\right|$.
	 \shortans{8}
	 \loigiai{
	\immini{ 
		Ta có  $\overrightarrow{AB}+\overrightarrow{AC}+\overrightarrow{AD}=(\overrightarrow{AB}+\overrightarrow{AD})+\overrightarrow{AC}=\overrightarrow{AC}+\overrightarrow{AC}$.\\
		Gọi $E$ đối xứng với $A$  qua $C$, suy ra $\overrightarrow{AC}=\overrightarrow{CE}$.\\
		 Khi đó $\overrightarrow{AB}+\overrightarrow{AC}+\overrightarrow{AD}=\overrightarrow{AC}+\overrightarrow{AC}=\overrightarrow{AC}+\overrightarrow{CE}=\overrightarrow{AE}$.\\
		Ta có $AE=2AC=2\cdot  2\sqrt{2}\cdot\sqrt{2}=8$.\\
		Vậy  $\left|\overrightarrow{AB}+\overrightarrow{AC}+\overrightarrow{AD}\right|=AE=8$.
		}{\begin{tikzpicture}[line join=round,line cap=round,>=stealth,scale=.6,font=\footnotesize]
		\foreach \x/\y/\n in  {0/0/A,0/2/B,2/0/D} \coordinate (\n) at (\x,\y);
		\coordinate (C) at  ($(B)+(D)-(A)$);
		\coordinate  (E) at ($(A)!2!(C)$);
		\draw  (A)--(B)--(C)--(D)--(A)--(E);
		\foreach  \t/\g in {A/-135,C/0,B/135,D/-45,E/45}{
			\draw[fill=black] (\t) circle (1pt)  node[shift={(\g:8pt)}]{$ \t $};
		} 
	\end{tikzpicture}
		} 
	} 
\end{bt}


%%%%%%%%%%%%%%%% TỰ LUẬN %%%%%%%%%%%%%%%%

% nguồn: MaHoaNH_TuLuan_Lop10
\begin{bt}%[0H5H2-4]%[ma:03e712d6]
	Cho  tam giác $ABC$ vuông tại $A$ có $\widehat{ABC}=30^{\circ}$ và $BC=\sqrt{5}$.
	Tính độ dài của véc-tơ $\overrightarrow{AB}+\overrightarrow{AC}$. 
	\loigiai{ 
	\immini{Gọi $D$ là điểm sao cho tứ giác $ABDC$ là hình bình  hành.\\
		Khi đó theo quy tắc hình bình  hành ta có $\overrightarrow{AB}+\overrightarrow{AC}=\overrightarrow{AD}$.\\
		Vì $\triangle ABC$ vuông ở $A$ nên tứ giác $ABDC$ là hình chữ nhật.\\ 
		Suy  ra $AD=BC=a\sqrt{5}$.\\
		Vậy $\left|\overrightarrow{AB}+\overrightarrow{AC}\right|=\left|\overrightarrow{AD}\right|=AD=a\sqrt{5}$. 
		}{\begin{tikzpicture}[line join=round,line cap=round,>=stealth,scale=1,font=\footnotesize]
		\foreach \x/\y/\n in {0/0/A,3/0/B,0/1.5/C}  \coordinate (\n) at (\x,\y);
		\coordinate (D) at ($(B)+(C)-(A)$); 
		\draw  (A)--(B)--(C)--(D)--(A)--(C)
			(B)--(D); 
		\foreach \t/\g in {A/-135,D/45,C/135,B/-45}{ 
			\draw[fill=black] (\t) circle (1pt) node[shift={(\g:8pt)}]{$  \t $};
		} 
	\end{tikzpicture}
		 }
	 }
\end{bt}

% nguồn: MaHoaNH_TuLuan_Lop10
\begin{bt}%[0H5H2-4]%[ma:5a13497a]
	Cho tam giác $ABC$ đều có  cạnh bằng $3 a$. Tính các giá trị sau
	\begin{enumerate}
		\item  $\left|\overrightarrow{AB}-\overrightarrow{AC}\right|$.
		\item  $\left|\overrightarrow{AB}-\overrightarrow{CB}-\overrightarrow{AC}\right|$.
		 \item $\left|\overrightarrow{AB}+\overrightarrow{AC}\right|$.
	\end{enumerate}
	\loigiai{ 
		\begin{enumerate}
			\item $\left|\overrightarrow{A  B}-\overrightarrow{A C}\right|=|\overrightarrow{CB}|=CB=3a$.
			\item $\left|\overrightarrow{AB}-\overrightarrow{CB}-\overrightarrow{AC}\right|=\left|\overrightarrow{AB}-\overrightarrow{AC}-\overrightarrow{CB}\right|=\left|\overrightarrow{CB}-\overrightarrow{CB}\right|=\left|\overrightarrow{0}\right|=0$. 
			\item Dựng hình bình hành  $ABDC$, ta có $\overrightarrow{AB}+\overrightarrow{AC}=\overrightarrow{AD}$.
		\immini{ 
			Gọi $O$ là giao điểm của $AD$ và $BC$.  \\
			 Ta có $AO=3 a\dfrac{\sqrt{3}}{2}$.\\
			 Suy ra $\left|\overrightarrow{AB}+\overrightarrow{AC}\right|=\left|\overrightarrow{AD}\right|=AD=2AO=3a\sqrt{3}$.
			}{\begin{tikzpicture}[line join=round,line cap=round,>=stealth,scale=1,font=\footnotesize]
		 \foreach \x/\y/\n in {0/0/B,2/0/C} \coordinate (\n) at (\x,\y);
		\coordinate (A) at ($(B)+({2*cos (60)},{2*sin  (60)})$);
		\coordinate (D)  at ($(B)+(C)-(A)$);
		\coordinate (O)  at ($(B)!0.5!(C)$);
		\draw  (A)--(B)--(C)--(D)--(A)--(C)
			 (B)--(D);
		 \foreach \t/\g in {A/90,D/-90,C/0,B/180,O/-45}{
			\draw[fill=black]  (\t) circle (1pt) node[shift={(\g:8pt)}]{$ \t $};
		} 
		\draw pic[thin,draw,angle radius=2mm] {right angle=  A--O--B};
	\end{tikzpicture}
			 }
		\end{enumerate}
	 }
\end{bt}

% nguồn: MaHoaNH_TuLuan_Lop10
\begin{bt}%[0H5H2-4]%[ma:c3cc1f3a]
	Cho hình vuông  $ABCD$ có cạnh bằng $a$ với tâm là $O$. Tính
	\begin{enumEX}{2}
		\item Độ dài  vectơ $\overrightarrow{OA}-\overrightarrow{CB}$.
		 \item Tính $\left|\overrightarrow{AB}+\overrightarrow{DC}\right|$.
	\end{enumEX}
	\loigiai{ 
		\begin{center}
			\begin{tikzpicture}[scale=1,font=\footnotesize,line join=round,line cap=round,>=stealth]
					 \def\a{3}
					\path (0:0) coordinate (A) 
						(0:\a)  coordinate (B)
						(90:\a)  coordinate (D)
						($(D)-(A)+(B)$) coordinate (C) 
						($2*(B)-(A)$)  coordinate (A')
						(intersection of  A--C and B--D) coordinate (O);
					\draw  (A)-- (B)-- (C)--(D)--cycle (A)-- (C)
						(B)-- (D) 
						(B)--(A'); 
					\foreach \x/\g  in {A/180,B/-90,C/45,D/155,O/-90,A'/0}
					\fill[black] (\x)  circle (1pt)
						($(\g:3mm)+(\x)$) node {$\x$}; 
				\end{tikzpicture}
		\end{center}
		\begin{enumerate}
			\item Ta có  $\overrightarrow{OA}-\overrightarrow{CB}=\overrightarrow{CO}-\overrightarrow{CB}=\overrightarrow{BO}$.\\
			Mặt khác  $BO=\dfrac{1}{2} BD=\dfrac{1}{2} \sqrt{a^2+a^2}=\dfrac{a\sqrt{2}}{2}$ nên $\left|\overrightarrow{OA}-\overrightarrow{CB}\right|=\dfrac{a\sqrt{2}}{2}$.
			 \item Gọi $A'$ là điểm đối xứng với $A$ qua $B$.
		\end{enumerate}
		Ta có $\overrightarrow{AB}+\overrightarrow{DC}=\overrightarrow{AB}+\overrightarrow{BA'}=\overrightarrow{AA'}$  nên $\left|\overrightarrow{AB}+\overrightarrow{DC}\right|=\left|\overrightarrow{AA'}\right|=2a$.
	 }
\end{bt}
